\section{Introduction}

\citet{hu2024layered} develop a layered modal type theory which supports pattern
matching on code. %
A critical idea lying in this system is \emph{the layering principle}. %
The layering principle begins with a core language, e.g. simply typed
$\lambda$-calculus (STLC). %
Then we extend the core language with a layer of the $\square$ modality which supports
meta-programming and intensional analysis. %
Though only two layers are demonstrated by \citet{hu2024layered}, in principle, this
extension can be iterated indefinite number of times, forming an arbitrary $n$-layered
modal type theory. %

In another dimension, instead of adding more and more layers, we could also increase
the expressive power of the core language. %
One interesting candidate for a core language is Martin-L\"of type theory (MLTT). %
MLTT is the foundation for many type-theory-based proof assistants, including Coq,
Agda and Lean. %
Treating MLTT as the core language and applying the layering principle to it could
yield a dependently typed system that allows to meta-program and intensionally analyze
code of itself without forgoing the consistency of the overall system. %
This feature gives a solid foundation for proof assistants to support truly type-safe
meta-programming. %
Due to the layering principle, libraries written for bare MLTT can also be used during
meta-programming. %
For example, we can use the same data structures for natural numbers and lists for
both programs and meta-programs. %
Meanwhile in reality, e.g. in Coq, we have at least four unexchangeable notions
natural numbers: the natural numbers defined inductively in Gallina, Ltac's natural
numbers, failure levels and hint database's search levels. %
Therefore, we can foresee that the layering principle also has the additional benefit
in engineering.

In this technical report, we first extend the previous layered modal type theory with
context variables. %
They are necessary to enable recursion on code. %
We justify the decidability of conversion checking following
\citet{abel_decidability_2017}'s reducibility proof. %
We then scale the setup all the way to MLTT, introducing \delamlang,
\textbf{De}pendent \textbf{La}yered \textbf{M}odal type theory. %
We then scale the reducibility proof to \delamlang and therefore justify its
decidability of conversion checking. %
A corollary is the consistency of \delamlang, hence showing that this type theory can
be used as a foundation for proof assistants.
